\section{Experiments}
\label{sec:experiments}

The experiments test whether aligning the training objective with ranked detections improves event localization after the detector components around the loss are controlled.
All methods are trained as sequence predictors and evaluated as event detectors.
BDL ranks local maxima of its learned score sequences.
Segmentation baselines estimate samplewise interval labels and then score transitions in that prediction.
Within each comparison we hold the series-level splits, backbone, scorer, and post-processing search fixed whenever possible, so the comparison isolates the supervision rather than the matching protocol.

The main experiments use two benchmarks.
The Child Mind Institute sleep benchmark contains wrist-accelerometer recordings annotated with sleep onset and wake-up events \citep{alexander2017,migueles2019,robbins2019,esper2023}.
It highlights the supervision mismatch because the annotated sleep interval can last for hours while evaluation scores minute-scale onset and wake-up matches.
The CHB-MIT seizure benchmark contains EEG recordings with seizure intervals sampled at 256 Hz and tests rarer events with second-scale matching tolerances \citep{shoeb2010,goldberger2000,klem1999}.
In both cases, supervision is built only for the event types scored by average precision: onset and wake-up for sleep, and seizure onset and offset for CHB-MIT.
The two benchmarks are deliberately different.
Sleep makes the mismatch visible because a long state interval is evaluated through short endpoint matches; CHB-MIT stresses the same formulation under higher sampling rate, rarer events, and stricter absolute timing.
\IfFileExists{results/generated/point_event_ablation_ready.tex}{Additional point-event ablations on Martian bow-shock and credit-card fraud detection test the same objective families without paired interval endpoints \citep{azib2023universal}.}{}

The sleep study first isolates the objective before varying the encoder.
With a bidirectional GRU fixed, we compare BDL-Hard, BDL-Gaussian, and BDL-Tolerance against cross-entropy segmentation.
We add class-weighted cross-entropy and focal segmentation to test whether standard imbalance handling explains the result \citep{krawczyk2016,lin2017focal}.
We then repeat the strongest BDL and segmentation objectives with one-dimensional U-Nets, an attention-gated U-Net variant, and an offline Transformer candidate \citep{ronneberger2015,oktay2018attention,vaswani2017attention}.
The U-Net controls test whether the same objective effect appears after replacing the recurrent encoder with convolutional segmentation backbones whose multi-scale skip connections are designed to recover time-localized labels from long context.
The attention-gated variant remains within this convolutional family: it changes skip-feature selection without changing the detector into a self-attention model.
If cross-entropy were weak only because the GRU was a poor interval-mask predictor, these segmentation-oriented backbones should narrow the gap before any specialized event decoder is introduced.
The Transformer comparison adds a self-attention detector family with different sequence-length, memory, and optimizer tradeoffs.
Its objective comparison still uses matched scoring and decoder calibration, but the result is interpreted as backbone-specific evidence rather than as a simple skip-connection variant.
Optimization settings are selected at the benchmark-and-backbone level for stable multi-epoch training, then held fixed inside each matched objective comparison.
This keeps learning rate, memory, and gradient-stability choices from becoming objective-specific advantages.
For sleep, the shared learning rate is chosen for reliable Poisson-objective training before the matched objective comparison and then reused for the segmentation controls.
For CHB-MIT, the high-capacity GRU configurations use a smaller step size, batch size 8, and gradient clipping because the EEG sequences are longer and the recurrent state is larger.
Calibration and post-processing parameters, such as candidate smoothing, score cutoffs, peak separation, and interval alternation, are selected on validation folds before aggregation.
They are therefore part of the evaluated detector, but not part of the training objective.
This protocol treats optimization and calibration as controlled experimental conditions rather than hidden method components.
The candidate-smoothing grid is chosen to match each benchmark's scoring scale.
For the reported CHB-MIT tables, decoder calibration is performed after expanding predictions back to the original 256 Hz timeline, with no smoothing, one-second smoothing, and two-second smoothing as the candidate scales.
This post-training rescoring keeps calibration aligned with second-scale seizure endpoint tolerances and avoids treating broad sleep-scale smoothing as a default for EEG endpoint scoring.

Each ablation is tied to a specific alternative explanation.
Kernel variants, sparse-output initialization, and smoothing-width comparisons ask whether the result depends on a particular smoothing or calibration convention inside BDL.
Because every BDL kernel assigns one event to each annotation before binning, these comparisons change the uncertainty model, initialization, or output scale without changing the event quantity being fit.
Class-weighted and focal segmentation keep the interval-mask labels and ask whether ordinary class-imbalance handling is enough.
Architecture controls ask whether the conclusion is specific to a recurrent encoder.
CHB-MIT changes the signal domain and timing scale; point-event experiments remove paired endpoints; output-stride studies expose the tradeoff among localization resolution, event sparsity, sequence length, and optimization; streaming-context experiments remove future samples from the encoder; and state reconstruction changes the evaluation object back to samplewise labels.
Together, these comparisons ask whether the improvement follows from event-level supervision or from model, calibration, and evaluation choices.
An advantage that survives imbalance controls, encoder changes, and event-structure changes supports the training-objective explanation, while drops under stricter stride, context, or calibration settings identify detector limitations around the likelihood rather than contradictions in the event formulation.

The BDL variants differ only in the unit-sum kernel used before temporal binning.
BDL-Hard places the event at the annotated timestamp.
BDL-Gaussian spreads it with a truncated Gaussian label-uncertainty model.
BDL-Tolerance uses a unit-sum kernel whose support is derived from the benchmark tolerance set.
All three are trained with the same Poisson log score.
The segmentation variants keep the samplewise interval-state labels and vary only the class-imbalance treatment.

Evaluation follows the benchmark output format.
Candidate smoothing, peak selection, and interval alternation are selected on validation folds.
These post-processing choices are held outside the training objective: they define how a trained score sequence is converted into ranked detections.
Sleep AP averages onset and wake-up over one-to-thirty-minute tolerances; seizure AP uses one-to-sixty-second tolerances, with one-to-three-second AP treated as the strict boundary-localization regime.
The output stride is selected with the metric-resolution, event-sparsity, and sequence-length heuristic in Appendix~\ref{app:derivation}.
For high-rate EEG, a 256-sample stride corresponds to one second at the original sampling rate, matching the strictest endpoint tolerance scale while keeping recurrent training feasible.
Coarser two-second settings are still reported as compute--accuracy controls, because output stride changes localization quantization, recurrent sequence length, and optimization difficulty at the same time.
\IfFileExists{results/generated/seizure_highscore_ready.tex}{The high-capacity CHB-MIT comparison evaluates a coarser two-second output stride against a matched one-second stride, separating the training objective from output-timeline resolution.}{}
We report validation-fold event mAP over series-level splits and reserve samplewise state-label metrics for the appendix analysis.
